\documentclass[11pt]{article}
\usepackage[margin=1in]{geometry}
\usepackage{graphicx,booktabs,amsmath,amssymb}
\usepackage{xcolor}
\usepackage[colorlinks=true,linkcolor=blue,citecolor=blue,urlcolor=blue]{hyperref}
\usepackage{natbib}
\usepackage{authblk}

\title{\textbf{Internals Know More Than the Model Can Say:\\
Relevance Signals Are Encoded Internally but Attenuated en Route to the Output}}
\author{Chenhui Feng}
\affil{Nanjing University of Aeronautics and Astronautics}
\date{}

\newcommand{\auc}{\textsc{auc}}

\begin{document}
\maketitle

\begin{abstract}
We study whether large language models (LLMs) ``know'' more about retrieval
relevance than their outputs reveal. Using linear probes on a single
frozen forward pass over (query, passage) pairs, we compare four
relevance \emph{readers} on the same model, prompt, and answer token:
(i) the model's literal output verdict (zero-shot LLM-judge), (ii) a trained
probe on the final-layer hidden state (the \emph{output representation}),
(iii) a trained probe on the best residual-stream layer (\emph{internal}),
and (iv) a probe on the top attention heads. Across eight models (3B--8B,
LLaMA, Mistral and Qwen, base and instruct), reading internal activations decodes
relevance markedly better than reading the output. With query-clustered
bootstrap confidence intervals, the attention-head probe beats the
output-representation probe on all four LLaMA models (each
$p<10^{-3}$ at the 1500-query scale) and on both Qwen models
($p<0.002$); the
residual-stream advantage over the output representation is individually
significant on base models ($+0.077$, $p{=}0.0008$; $+0.060$, $p{=}0.004$;
Qwen $+0.063$, $p{=}0.001$).
A causal steering intervention confirms the link is mechanistic, not
epiphenomenal: ablating the internal relevance direction at the output layer
monotonically collapses the model's verdict (mean $P(\text{yes})$
$0.79{\to}0.26$) while keeping the verdict discriminative (\auc{} $\approx
0.55$ across the sweep). It is the only direction that does both: the output's
own direction shifts $P(\text{yes})$ but collapses \auc{} to $0.485$ (a
non-discriminative logit shift), and a random direction leaves \auc{} flat
without a relevance-aligned verdict shift. Comparing base and instruct pairs, we
find alignment improves the output readout (the judge gets stronger) yet does
not close the internal--output gap. In short: the relevance signal is
encoded internally and \emph{attenuated}, not absent, on the path to the
logits.
\end{abstract}

\section{Introduction}
When an LLM is asked whether a retrieved passage is relevant to a query, its
literal answer is often little better than chance. Yet the same forward pass
may already contain a far cleaner relevance signal in its internal
activations. We ask a focused, mechanistic question: \emph{is relevance
better decodable from a model's internals than from its output, and if so,
is the internal signal causally connected to the output the model fails to
produce?}

This is an interpretability question, not a retrieval-engineering one. We do
\textbf{not} claim a probe that beats a strong cross-encoder reranker, nor
zero-shot cross-domain robustness; prior exploration refuted those framings.
Our contribution is the \emph{relative} science of internals vs.\ output:

\begin{itemize}
  \item \textbf{Existence.} On the same model/prompt/answer-token, a linear
  probe on attention-head activations decodes (query, passage) relevance
  substantially better than the model's own output verdict, and better than a
  trained probe on the output representation (final hidden state).
  \item \textbf{Attribution.} Decomposing the gap into a
  \emph{training} component (final-layer probe $-$ judge) and a
  \emph{depth/internal} component (best internal layer $-$ final layer) shows
  the edge comes substantially from \emph{reading internals}, not merely from
  \emph{being trained}.
  \item \textbf{Mechanism.} Relevance decodability peaks in the middle of the
  network and declines toward the output layer. Comparing base vs.\ instruct
  models, alignment \emph{improves} the output readout but does not eliminate
  the internal--output gap.
  \item \textbf{Causality.} A steering intervention that injects the internal
  relevance direction at the output layer causally controls the verdict;
  ablation is cleanly monotonic and, unlike the control directions, both moves
  the verdict and keeps it discriminative---evidence that the output pathway
  \emph{can} express the signal but normally under-expresses it.
\end{itemize}

We report query-clustered bootstrap confidence intervals throughout, and are
explicit about which claims are individually significant, which rest on a
cross-model sign test, and which are architecture-dependent (the
attention-head reader is weak on Mistral). Honest delineation of effect
strength is central to this paper.

\section{Related Work}
\textbf{Probing classifiers} read latent features from frozen representations
to test what information is linearly available \citep{alain2017,hewitt2019};
we use them comparatively across depth and against the output.
\textbf{Logit lens and tuned lens} \citep{nostalgebraist2020,belrose2023}
project intermediate states to the vocabulary, motivating our depth analysis.
\textbf{Activation steering / representation engineering}
\citep{turner2023,zou2023,li2023iti} edit residual directions to change
behavior; we use a difference-of-means steering vector for a causal test.
\textbf{Calibration and hallucination} work shows models can encode
correctness they do not surface \citep{kadavath2022,azaria2023,burns2022}.
Most relevant, \citet{orgad2024} show LLMs internally represent truthfulness
they do not output, and concurrent work finds hallucination is linearly
decodable from mid-layer hidden states while output-based detectors stay near
chance \citep{aiersilan2026}. We do \emph{not} claim to be first to observe
that internals outstrip the output; our contribution is to localize this to a
specific, controllable signal---query--passage \emph{relevance}---and to
characterize it quantitatively: a depth-resolved attenuation curve, a
decomposition separating ``reading internals'' from ``being trained,'' a causal
steering test, and a base-vs-instruct contrast. \textbf{Relevance mechanism}
work reverse-engineers how neural rankers compute relevance via causal
interventions on attention heads \citep{chen2024relevance}; trained classifiers
are also known to beat prompted LLM judges at relevance \citep{gienapp2025}. We
differ by studying a generative LLM's internal-vs-output gap on relevance rather
than a ranker's internal circuit, and by attributing the gap to depth beyond
training. \textbf{RAG relevance} estimation usually targets reranking accuracy;
we deliberately stay on the interpretability question and do not compete with
rerankers.

\section{Method}
\label{sec:method}
\textbf{Task and data.} We use MS MARCO v1.1, forming (query, passage, label)
triples with relevance labels. Splits are made \emph{by query} (70/15/15) so
no query appears in both train and test, preventing leakage. We report
ROC-\auc{} on the held-out test set.

\textbf{Prompt and extraction.} Every reader sees the identical instruct
prompt, \texttt{"Document: \{p\}\textbackslash nQuestion:
\{q\}\textbackslash nIs the document relevant to the question? Answer:"}, and
all activations are taken at the final \emph{answer token}. Per-head
activations are the \texttt{o\_proj} outputs reshaped to
$[\text{heads}, d_{\text{head}}]$; the residual stream is each decoder
layer's output.

\textbf{Four readers.} (1) \emph{LLM-judge}: the model's own normalized
$P(\text{yes})$ vs.\ $P(\text{no})$ at the answer token, \emph{zero training}.
(2) \emph{Output-representation probe}: an $\ell_2$-regularized logistic
regression on the final hidden state. (3) \emph{Internal probe}: the same
probe on the best residual-stream layer, selected \emph{on validation}.
(4) \emph{Attention-head probe}: top-$k$ heads chosen on validation, their
activations concatenated into one logistic regression. As an untrained
output-side control we also use a \emph{logit lens}: each layer's residual is
projected through the model's own final norm and unembedding to read
$P(\text{yes})$ vs.\ $P(\text{no})$ at every depth, with no training and no
label-based layer selection.

\textbf{Gap decomposition.} We split the total edge over the judge into
$\text{train\_edge} = \text{final} - \text{judge}$ (pure training, both at the
output) and $\text{internal\_edge} = \text{best\_resid} - \text{final}$ (pure
depth, beyond training). The thesis ``internals $>$ output'' requires internal
readers to exceed \emph{both} the judge and the output-representation probe.

\textbf{Significance.} Because the split is by query, the correct resampling
unit is the query. We report 95\% confidence intervals and one-sided
$p$-values from a \emph{query-clustered bootstrap} (5000 resamples) on frozen
test predictions; the best layer is chosen on validation to avoid
selection-on-test optimism. We additionally run a cross-model sign test on
the direction of each edge.

\textbf{Causal steering.} At the best internal layer we form a probe-free
relevance direction $d = \overline{h}_{\text{rel}} - \overline{h}_{\text{irrel}}$
(difference of class means, unit-normalized). During the forward pass we add
$\alpha \|h\| d$ to the \emph{output-layer} residual at the answer token and
sweep $\alpha$, reading the judge's $P(\text{yes})$ and the test \auc. We
compare against the output layer's own difference-of-means direction and a
random unit direction (norm-matched).

\section{Experiments}
\label{sec:exp}

\textbf{Internals decode relevance better than the output.}
Table~\ref{tab:main} reports the four readers across eight models. On every
model the best internal reader exceeds the trained output-representation probe,
and on every model except the Qwen pair it also exceeds the zero-shot judge.
The judge hovers near chance on the LLaMA and Mistral base models (\auc{}
$0.51$--$0.54$), confirming the output ``says'' little; Qwen is the informative
exception---its zero-shot judge is already strong ($0.65$), yet reading its
attention heads still beats reading its output representation (below).

\begin{table}[t]
\centering\small
\caption{Four relevance readers across eight models (test \auc, point
estimates; val-selected layers and heads). On every model the attn probe
exceeds both the judge and the output-representation probe; best-resid also
exceeds both on the LLaMA and Mistral models. Qwen's zero-shot judge is
notably stronger ($0.65$) than the LLaMA/Mistral base models ($0.51$--$0.54$);
despite this, reading attention heads still beats the output representation on
both Qwen models (see Table~\ref{tab:scale}).}
\label{tab:main}
\begin{tabular}{lcccccc}
\toprule
Model & $L$ & judge & final & best-resid & attn \\
\midrule
LLaMA-3.2-3B (base)      & 28 & 0.535 & 0.592 & 0.628 & \textbf{0.685} \\
LLaMA-3.1-8B (base)      & 32 & 0.514 & 0.552 & 0.581 & \textbf{0.680} \\
Mistral-7B-v0.3 (base)   & 32 & 0.531 & 0.601 & \textbf{0.621} & 0.603 \\
LLaMA-3.2-3B-Instruct    & 28 & 0.597 & 0.606 & 0.633 & \textbf{0.725} \\
LLaMA-3.1-8B-Instruct    & 32 & 0.603 & 0.583 & 0.628 & \textbf{0.746} \\
Mistral-7B-Instruct-v0.3 & 32 & 0.621 & 0.616 & 0.628 & \textbf{0.667} \\
Qwen2.5-7B (base)        & 28 & 0.652 & 0.592 & 0.655 & \textbf{0.682} \\
Qwen2.5-7B-Instruct      & 28 & 0.660 & 0.610 & 0.630 & \textbf{0.670} \\
\bottomrule
\end{tabular}
\end{table}

\textbf{Significance.} Table~\ref{tab:sig} gives query-clustered bootstrap
CIs. The \emph{attention-head vs.\ output-representation} edge is individually
significant on all four LLaMA models (effect $+0.09$ to $+0.16$): $p\le0.016$
at 500 queries (Table~\ref{tab:sig}) and $p<10^{-3}$ once the evaluation is
scaled (Table~\ref{tab:scale}).
The residual-stream edge (internal\_edge) is directionally positive on all six
models in the 500-query evaluation (sign test $p{=}0.016$) but, at $\sim$75 test
queries, not individually significant---a power problem, not an absent effect
(confirmed by scaling, below).

\begin{table}[t]
\centering\small
\caption{Query-clustered bootstrap (5000 resamples), 500-query evaluation.
Edges with 95\% CI and one-sided $p$. \textbf{Bold}: $p<0.05$.}
\label{tab:sig}
\begin{tabular}{lccc}
\toprule
Model & internal\_edge & attn$-$final & best\_resid$-$judge \\
\midrule
LLaMA-3.2-3B      & $+0.036$ [-.05,.12] & \textbf{$+0.094$} [.01,.18] & \textbf{$+0.093$} [.01,.18] \\
LLaMA-3.1-8B      & $+0.028$ [-.06,.11] & \textbf{$+0.128$} [.04,.22] & $+0.065$ [-.02,.15] \\
Mistral-7B        & $+0.019$ [-.07,.10] & $+0.002$ [-.08,.09]        & \textbf{$+0.089$} [.00,.17] \\
LLaMA-3.2-3B-Inst & $+0.026$ [-.06,.11] & \textbf{$+0.119$} [.03,.21] & $+0.035$ [-.05,.13] \\
LLaMA-3.1-8B-Inst & $+0.044$ [-.04,.13] & \textbf{$+0.164$} [.08,.25] & $+0.024$ [-.06,.11] \\
Mistral-7B-Inst   & $+0.011$ [-.08,.10] & $+0.051$ [-.04,.15]        & $+0.006$ [-.08,.09] \\
\midrule
sign test ($p$)   & 6/6 (\textbf{.016}) & 6/6 (\textbf{.016}) & 6/6 (\textbf{.016}) \\
\bottomrule
\end{tabular}
\end{table}

\textbf{Scaling the evaluation confirms the internal edge.}
We re-extract four flagship models at 1500 queries ($\sim$225 test queries;
sample size is the only change) and re-test (Table~\ref{tab:scale}). The
residual-stream edge becomes individually significant on base models
($+0.077$, $p{=}0.0008$; $+0.060$, $p{=}0.004$): the $p$-values drop from the
$0.2$ range to the $0.001$ range purely from added power. Strikingly, the edge
is significant on base models but shrinks and loses significance on instruct
models---a mechanistic signal we return to below.

\begin{table}[t]
\centering\small
\setlength{\tabcolsep}{4.5pt}
\caption{Scaled (1500-query) re-test vs.\ the 500-query estimate. The
internal\_edge becomes individually significant on base models once power is
added; the attn$-$final edge is significant on all four LLaMA models at this
scale. \textbf{Bold}: $p<0.05$.}
\label{tab:scale}
\begin{tabular}{lccc}
\toprule
Model & internal\_edge (500q) & internal\_edge (1500q) & attn$-$final (1500q) \\
\midrule
LLaMA-3.2-3B (base) & $+0.036$, $p{=}0.19$ & \textbf{$+0.077$ [.03,.12], $p{=}0.0008$} & \textbf{$+0.146$ [.11,.19], $p{<}10^{-3}$} \\
LLaMA-3.1-8B (base) & $+0.028$, $p{=}0.26$ & \textbf{$+0.060$ [.02,.10], $p{=}0.004$} & \textbf{$+0.114$ [.08,.15], $p{<}10^{-3}$} \\
LLaMA-3.2-3B-Inst   & $+0.026$, $p{=}0.28$ & $+0.024$ [-.02,.07], $p{=}0.15$ & \textbf{$+0.097$ [.06,.13], $p{<}10^{-3}$} \\
LLaMA-3.1-8B-Inst   & $+0.044$, $p{=}0.16$ & $+0.012$ [-.04,.06], $p{=}0.32$ & \textbf{$+0.129$ [.09,.17], $p{<}10^{-3}$} \\
Qwen2.5-7B (base)   & --- & \textbf{$+0.063$ [.02,.10], $p{=}0.001$} & \textbf{$+0.090$ [.05,.13], $p{<}10^{-3}$} \\
Qwen2.5-7B-Inst     & --- & $+0.020$ [-.01,.05], $p{=}0.12$ & \textbf{$+0.060$ [.02,.10], $p{=}0.002$} \\
\bottomrule
\end{tabular}
\end{table}

\textbf{A third model family (Qwen2.5) confirms and extends the pattern.}
Qwen2.5-7B has a notably stronger zero-shot judge than LLaMA or Mistral
($0.65$ vs.\ $0.51$--$0.54$), suggesting its pretraining already embeds
relevance signals more directly in the output vocabulary direction. Despite
this, reading the attention heads still significantly beats the output
representation on both Qwen models (base $+0.090$, $p{<}10^{-3}$; instruct
$+0.060$, $p{=}0.002$). The internal\_edge is significant on the base model
($+0.063$, $p{=}0.001$) and shrinks on the instruct model---the same
base-to-instruct compression seen on LLaMA, replicated in a third family.
The train\_edge is \emph{negative} on both Qwen models ($-0.060$,
$-0.050$): once the native judge is already strong, training a probe on the
final hidden state actually slightly underperforms it. This is the structural
mirror of the LLaMA case---where alignment \emph{moved} signal toward the
output (judge rising from $0.52$ to $0.60$), Qwen appears to have that
signal \emph{already present} in the output pathway from pretraining. The
internal--output gap persists regardless.

\textbf{Depth localization.} Figure~\ref{fig:depth} (per model in the
supplement) shows relevance decodability rising into the mid network and
declining toward the output layer; the judge baseline sits well below the
peak. For LLaMA the strongest heads cluster in the middle (L10--16); for
Mistral single heads carry little signal (early-layer head \auc{} $\approx
0.50$) while the residual stream still encodes relevance---a real
architectural difference, not a failure of the thesis.

\textbf{The signal is not surfaced by the model's own unembedding (logit
lens).} One might object that mid-network decodability only reflects what a
\emph{trained} probe can extract, not what the model could surface. We test
this with a logit lens: at each layer we project the cached residual through
the model's own final norm and unembedding and read the normalized
$P(\text{yes})$ vs.\ $P(\text{no})$---an untrained, output-side reader at every
depth. As a sanity check, the last-layer logit lens reproduces the judge
exactly (3B and 8B: $0.545$ and $0.574$, matching the judge to within
$<0.001$). Crucially, the logit lens stays near chance at \emph{every} depth
(3B peak $0.560$ at L25; 8B peak $0.577$ at L30), far below the trained
residual probe's mid-network $0.62$. So the relevance signal present in the
residual stream is not readable through the model's vocabulary head at any
layer: the bottleneck is not that the signal is absent from late layers but
that it is \emph{misaligned with the output readout}, consistent with the
causal result that the output's own direction shifts the verdict without
preserving discriminative \auc.

\begin{figure}[t]
\centering
\includegraphics[width=0.92\textwidth]{figures/llama32_3b_layer_signal.png}
\caption{Depth localization (LLaMA-3.2-3B). (a) per-layer residual probe \auc{}
vs.\ depth with the judge baseline and the peak/output markers; (b) per-head
\auc{} heatmap; (c) head- vs.\ layer-view and the layer distribution of the top
heads. Signal forms mid-network and attenuates toward the output.}
\label{fig:depth}
\end{figure}

\textbf{Mechanism: base vs.\ instruct.} Across three base$\to$instruct pairs,
alignment shows a consistent triple effect: the judge improves (output
``says'' more), train\_edge falls toward or below zero (training a final-layer
probe adds little once aligned), and the internal edge persists (and widens on
LLaMA). Combined with Table~\ref{tab:scale} (the internal edge is significant
on base, compressed on instruct), the reading is: \emph{alignment improves the
output readout but does not close the internal--output gap}---it does not
merely ``suppress'' the output.

\textbf{Causality.} Figure~\ref{fig:causal} shows the steering sweep. Injecting
the internal relevance direction at the output layer drives the verdict
monotonically on the ablation arm ($P(\text{yes})$ $0.79\to0.26$ as
$\alpha{<}0$) while keeping the verdict discriminative (\auc{} $0.546\to0.541$,
essentially flat). The three directions dissociate cleanly: the internal
direction both moves the verdict \emph{and} preserves \auc; the output's own
direction moves $P(\text{yes})$ strongly ($0.01\to0.99$) but \emph{destroys}
\auc{} ($\to0.485$, a uniform logit shift); and a random direction leaves \auc{}
flat ($\approx0.545$) but produces no relevance-aligned verdict shift. Only the
internal direction does both---the diagnostic signature of a causal,
relevance-bearing direction rather than a generic bias knob.
The additive arm ($\alpha{>}0$) saturates because the judge is already
yes-biased (baseline $0.79$ while only $\sim$22\% of test pairs are relevant),
so we base the causal claim on the ablation arm and the \auc{} contrast. This
indicates the output pathway \emph{can} express the internal relevance signal
but normally under-expresses it: attenuation, not absence. The same protocol
replicates across three variants (LLaMA-3.2-3B, LLaMA-3.1-8B base and
instruct): the ablation arm collapses $P(\text{yes})$ in every case, and only
the internal direction preserves \auc. The collapse is sharpest on the instruct
model ($P(\text{yes})$ $0.87\to0.003$), mirroring the alignment mechanism above:
alignment couples the verdict more tightly to the internal direction without
letting the output representation overtake the residual stream.

\begin{figure}[t]
\centering
\includegraphics[width=0.92\textwidth]{figures/causal_steer_3b.png}
\caption{Causal steering at the output layer (LLaMA-3.2-3B). (a) mean judge
$P(\text{yes})$ vs.\ steering coefficient $\alpha$; (b) test \auc. Only the
internal direction preserves \auc; the final-layer direction inflates
$P(\text{yes})$ while collapsing \auc.}
\label{fig:causal}
\end{figure}

\textbf{A second in-domain dataset.} To check the gap is not specific to
MS MARCO, we retrain and re-test the readers in-domain on BeIR FiQA
(LLaMA-3.2-3B, 1500 queries, 225 test queries; \emph{not} cross-domain
transfer). Negatives are sampled from corpus passages not judged relevant to
the query, so FiQA negatives are easier than MS MARCO's in-passage hard
negatives; this makes FiQA highly linearly separable and is why its
internal\_edge is thin (see below). The headline gap is even larger: the judge reaches only
$\auc{}=0.739$ while the best internal layer reaches $0.993$
(best\_resid$-$judge $=+0.253$, CI $[0.224,0.283]$). FiQA relevance is highly
linearly separable, so the output representation itself is near ceiling
($\auc{}=0.976$); the internal edge is correspondingly thin but still
significant ($+0.017$, CI $[0.010,0.024]$, $p<10^{-3}$), and depth still
peaks before the output (L14 $0.993$ vs.\ L28 $0.977$). The two datasets are
complementary --- MS MARCO stresses ``reading internals beats reading the
output representation,'' FiQA stresses ``a trained read-out far exceeds the
zero-shot output judge'' --- but both show the output judge trailing internal
decodability by a wide margin.

\textbf{A stronger judge does not close the gap.} A reviewer might attribute
the weak judge to a weak prompt. Adding four balanced in-context exemplars
(drawn only from train, no leakage) lifts the 3B base judge from
$\auc{}=0.545$ to only $0.566$ --- still far below the internal readers
(best-resid $0.614$, attn $0.682$ at the same 1500-query scale). The output
bottleneck is attenuation of the signal en route
to the logits (Fig.~\ref{fig:causal}), not prompt under-specification.

\section{Discussion and Limitations}
The relevance signal an LLM needs to judge a passage is present in its
internals and only partially surfaced at the output; alignment tightens the
readout without closing the gap. This connects to calibration/hallucination
findings that models ``know'' more than they assert, localized here to a
specific signal and resolved over depth, with a causal handle. The Qwen2.5
family is instructive: its zero-shot judge is already strong, so the
internals-vs-\emph{judge} gap nearly closes, yet reading attention heads still
beats reading the output representation---the internal--output gap is robust to
how well the output verbalizes relevance.

\textbf{Limitations.} (1) Probing is correlational; our causal evidence is a
steering study replicated on three LLaMA variants (3B, 8B base/instruct), whose
additive arm is limited by the judge's yes-bias, so we rely on the ablation arm.
(2) We use two in-domain relevance datasets (MS MARCO, FiQA); we make no
cross-domain or reranking-accuracy claims, which prior exploration refuted.
(3) The attention-head reader is architecture-dependent: it is the strongest
internal reader on LLaMA and Qwen but weak on Mistral (single heads near
chance), so we treat the residual-stream layer as the canonical ``internal''
reader and the head probe as a complementary instrument. (4) Absolute \auc{}
values are modest (0.5--0.75) on MS MARCO; the contribution is the \emph{relative}
internal-vs-output relationship, not a state-of-the-art relevance classifier.

\section{Conclusion}
On the same model, prompt, and token, relevance is more decodable from
internal activations than from the output the model produces, the advantage
comes substantially from reading depth rather than from training, the signal
peaks mid-network and attenuates toward the logits, and a steering ablation
shows the internal direction causally underlies the verdict. Internals know
more than the model can say.

\bibliographystyle{plainnat}
\bibliography{refs}

\end{document}
